%%%%%%%%%%%% ------------ chapter header ------------  %%%%%%%%%%%%
\chapterhead
% chapter number, title
{実験}
% chapter abst
{本章では，〜〜について述べる．}

%%%%%%%%%%%% ------------ next page ------------  %%%%%%%%%%%%

\section{実験環境}
実験を行った環境について説明する．
ハードウェア，ソフトウェア，ネットワーク環境などを記載する．

\section{実験条件}
実験の設定や条件を詳しく説明する．
パラメータ，データセット，被験者数などを明記する．

% 長い場合は\input{texファイル名}で挿入など工夫する

\section{評価指標}
実験結果をどのような指標で評価するのかを説明する．
精度，速度，ユーザビリティなど，具体的な評価基準を示す．

% 表の例
\begin{table}[h]
  \begin{center}
    \caption{システムAとシステムBの比較}
    \begin{tabular}{|l|l|l|} \hline
         システム　　 & A & B \\ \hline
         速さ 　　　　　　　　　& 1 & 2 \\ \hline
         手軽さ 　　　　　　& 3 & 4 \\ \hline
         持ちやすさ & 5 & 6 \\ \hline
         重さ　　　　　　　　　 & 7 & 8 \\ \hline
    \end{tabular}
    \label{tb:angle}
  \end{center}
\end{table}
% 参考：https://atatat.hatenablog.com/entry/cloud_latex7_table
% 位置
% []の中身
% h　figure環境宣言した場所（here）に図の場所を確保
% t　figure環境宣言したページの上部（top）に図の場所を確保
% b　figure環境宣言したページの下部（bottom）に図の場所を確保
% p　最後のページ（p）に図の場所を確保


% 数式
% $x = 1$　のように書く
\newpage

